\chapter{Thyroidectomy}
\section*{Overview}
Thyroidectomy removes part or all of the thyroid gland.
\section*{Common Instruments and Equipment}
The exact setup varies by technique and setting. Common items may include:
\begin{itemize}[nosep,leftmargin=*]
  \item neck surgery instrument set and magnification
  \item electrosurgical or vessel-sealing device
  \item nerve-monitoring equipment when indicated
\end{itemize}

\section*{Why It Is Done}
It may be performed for thyroid cancer, compressive goiter, suspicious nodules, or selected hyperthyroidism.
\section*{How It Is Performed}
Under general anesthesia, the surgeon removes one lobe or the entire gland through a lower-neck incision while protecting the recurrent laryngeal nerves and parathyroid glands.
\section*{Contraindications}
No single contraindication applies to every patient; the treating team assesses whether the procedure is suitable for the individual and the reason for treatment.

\section*{Precautions}
The care team reviews relevant health conditions, medicines, allergies, and procedure-specific risks, and explains any changes needed for the individual.

\section*{Risks and Recovery}
Bleeding, airway compromise, voice changes, hypocalcemia, infection, and need for lifelong thyroid hormone after total removal are possible.
\section*{References}
\begin{enumerate}\item MedlinePlus. Thyroidectomy. U.S. National Library of Medicine; 2025.\item Current Medical Diagnosis and Treatment. McGraw-Hill; 2025.\end{enumerate}
\clearpage
